\chapter{术语、符号与约定}
\section{对象与标度}

\begin{longtable}{@{}p{0.20\textwidth}p{0.34\textwidth}p{0.36\textwidth}@{}}
\caption{胶子 PDF/TMD 计算中最容易混淆的术语。}\label{tab:b10-terms}\\
\toprule
术语 & 定义 & 常见错误 \\
\midrule
\endfirsthead
\toprule
术语 & 定义 & 常见错误 \\
\midrule
\endhead
共线 PDF & 对横向动量积分后只保留纵向份额 $x$ 的分布；依赖 $\mu$。 & 把
quasi-PDF 直接当成共线 PDF，或把 TMD 的 $b_\perp$ 积分当成严格恒等式而不报告
幂修正。 \\
TMD-PDF & 同时依赖 $x,\bperp,\mu,\zeta$，含 staple 与软/快度结构。 & 只加一条
横向 Wilson 线而忽略软函数、过程方向和 CS 演化。 \\
quasi-PDF & 欧氏空间方向分离、由大 $\Pz$ 匹配到光锥 PDF 的准分布。 & 把有限
$\Pz$ 结果直接称为光锥结果。 \\
pseudo-ITD & 在 Ioffe 时间 $\nu=z\Pz$ 中表示同一矩阵元。 & 把坐标截断后的
振荡当成物理端点行为。 \\
RI/MOM & 在离壳夸克顶点和动量窗口内定义的非微扰 UV 方案。 & 对非局域 Wilson
线直接套用局域算符的减除，忽略 $z\mu$ 混合。 \\
hybrid & 短距 ratio 与长距自重整化/质量减除的拼接。 & 在 $z_s$ 两侧重复或遗漏
因子，或按点选择过渡尺度。 \\
CS 核 & 控制 TMD 对 $\zeta$ 演化的快度反常量纲。 & 把两动量比值的有效斜率
无条件外推到非微扰 $b_\perp$。 \\
SFTE & 正流时间算符对 $\MSbar$ 算符的小 $\tau$ 展开。 & 把流化后的 UV 有限性
误认为无需匹配，或把大流降噪结果用于小流外推。 \\
Wilson 圈减除 & 用同几何闭合圈吸收线性自能、pinch/cusp 结构。 & 用不同长度、
不同表示或不同转角几何的圈作减除。 \\
\bottomrule
\end{longtable}

\section{符号检查表}

\begin{tabularx}{\textwidth}{@{}p{0.15\textwidth}p{0.36\textwidth}X@{}}
\toprule
符号 & 含义 & 量纲/检查 \\
\midrule
$a$ & 格距 & fm 或 $\mathrm{GeV}^{-1}$；所有指数与对数必须同单位。 \\
$\tau$ & 流时间 & 物理量纲为长度平方；代码常用 $\tau/a^2$。 \\
$\sqrt{8\tau}$ & 流化半径 & 与 $|z|$、$b_\perp$ 比较时使用同一单位。 \\
$\Pz$ & 核子 boost 动量 & 不是 UV 标度 $\mu$，也不是快度标度 $\zeta$。 \\
$\nu=z\Pz$ & Ioffe 时间 & 无量纲；Fourier 共轭到 $x$。 \\
$Z_R$ & 自重整化因子 & 依赖 $z,a,\mu$；不能和硬匹配 $Z_{ij}$ 混为一谈。 \\
$K(\bperp,\mu)$ & Collins--Soper 核 & 依赖表示、Wilson 线几何和快度方案。 \\
$S_I(\bperp,\mu)$ & 内禀软函数 & 与完整软函数、$\zeta$ 方案分层记录。 \\
\bottomrule
\end{tabularx}

U055 的术语表适合作为附录索引，但它本身是综合性资料而非原始论文。其中弦张力
条目的数值与括号单位存在自相矛盾；专著应回到原始势能定义和自然单位换算，
不能原样复制该数字。U055 对 PDF、TMD、LaMET、RI/MOM、hybrid、CS 核和
Wilson 线仍可作为统一命名的有用底稿。

\chapter{来源导航、证据等级与版本谱系}

\section{六单元的使用方式}

\begin{longtable}{@{}p{0.13\textwidth}p{0.29\textwidth}p{0.48\textwidth}@{}}
\caption{B10 来源的证据定位。}\label{tab:b10-sources}\\
\toprule
单元 & 来源类型 & 专著中的用法 \\
\midrule
\endfirsthead
\toprule
单元 & 来源类型 & 专著中的用法 \\
\midrule
\endhead
U008 & 归档目录 README & 只证明旧 \texttt{examples/...} 路径和日志文件是历史快照；不能作为
现行运行入口或新物理证据。 \\
U033 & 2022 年准 PDF 学位论文 & 摘要级历史对照：RI/MOM 与 Wilson 线线性发散的
问题意识。正文数值与完整重整化条件未在本批次核对。 \\
U034 & 2023 年 TMD 学位论文 & 目录级工作流线索：TMD、Wilson loop、HMC/Chroma
和统计重采样。扫描 OCR 有字符损坏，未使用其公式。 \\
U055 & 格点 QCD 术语/缩写表 & 术语与符号附录的分类与索引；具体物理结论回溯到
U077 或原始文献，不把术语表当作独立推导。 \\
U078 & 旧版共线 PDF 工作流 & 历史版式与 Eq.(20)/Eq.(25) 叙事；已由 U077 取代，
只保留演进解释，不重复计入物理内容。 \\
U077 & 新版胶子 TMD 工作流 & 历史与复现章节的主要 TeX 证据：定义、SFTE、方案、
ratio-v2、匹配、验证边界和历史链。PDF 只用于页数与成对核对。 \\
\bottomrule
\end{longtable}

\section{行号、页码与 OCR 规则}

TeX/Markdown 来源使用“文件名:行号”锚定，例如 PDF/TMD 定义可回溯到
\texttt{理论解析与工作流-v20260911.tex:492--680}，算符与 TMD staple 可回溯到
\texttt{同文件:1220--1373}，SFTE 与窗口可回溯到 \texttt{同文件:2188--2300}。
这些行号指向固定快照，不随当前工作树的其他修改变化。

独立 PDF 使用物理 PDF 页码。U033 的第 2 页是中文摘要（印刷页 i），第 3 页
是英文摘要（印刷页 iii）；U034 的第 3 页是目录首页。扫描 PDF 的 OCR 出现
拆字、错码和字符重叠时，必须标为“公式/文本未可靠抽取”，不能根据残片补写
公式。PDF 页数与 TeX 自述不一致时，以来源文件的物理页数和 TeX 行为主记录，
并在 unresolved 中说明口径差异。

\section{证据等级}

\begin{enumerate}
  \item \textbf{A：定义或推导。} TeX 中给出对象、方程和适用条件的部分；
  仍需检查符号约定是否与专著主文一致。
  \item \textbf{B：来源内验证。} U077 登记的契约、回归、oracle 或参考诊断；
  未在 B10 重跑时只能写作“来源登记”。
  \item \textbf{C：摘要/目录级证据。} U033、U034 的题录、摘要和目录；可用于
  历史定位，不足以证明完整算法。
  \item \textbf{D：历史导航。} U008 的旧路径和 U078 的被替代版本；用于解释
  迁移与版本差异，不产生新的物理结论。
\end{enumerate}

任何数值进入正文前应回答四个问题：它来自哪个对象、在哪个尺度窗口、对应哪个
重整化/匹配方案、是否有独立复跑或只属于来源登记。若四个问题不能同时回答，
应把数值放入诊断表或未解决清单，而不是写成第一性原理结果。
